\providecommand{\Needspace}[1]{}
\Needspace{8\baselineskip}
\item \textbf{Movimiento de un proyectil lanzado desde un edificio}.
\nopagebreak[4]

\begin{tcolorbox}[colback=blue!2!white,colframe=blue!55!black,title={Enunciado},breakable]
Un proyectil se lanza desde la parte superior de un edificio de altura inicial $y_0$, con rapidez inicial $v_0$ y un ángulo de lanzamiento $\theta$ respecto de la horizontal. Se desea estudiar su movimiento usando arreglos de NumPy y visualizar, en una misma figura, la posición, velocidad y aceleración en las direcciones horizontal y vertical.

Considere los siguientes valores constantes:
\[
y_0 = 50 \ \text{m}, 
\qquad 
v_0 = 35 \ \text{m/s},
\qquad 
\theta = 55^\circ,
\qquad 
g = 9.8 \ \text{m/s}^2.
\]

El tiempo debe representarse mediante un arreglo temporal construido con \texttt{np.linspace}. En este problema use un arreglo \texttt{t} con $300$ valores entre $t=0$ y $t=8$ segundos.

Las ecuaciones cinemáticas del movimiento son:
\[
x(t) = v_0 \cos(\theta)t,
\]
\[
y(t) = y_0 + v_0 \sin(\theta)t - \frac{1}{2}gt^2.
\]

Las velocidades en cada dirección están dadas por:
\[
v_x(t) = v_0 \cos(\theta),
\]
\[
v_y(t) = v_0 \sin(\theta) - gt.
\]

Las aceleraciones correspondientes son:
\[
a_x(t) = 0,
\]
\[
a_y(t) = -g.
\]

Escriba un programa en Python que realice lo siguiente:

\begin{itemize}
    \item[(a)] Defina las constantes físicas del problema y construya el arreglo temporal $t$ usando \texttt{np.linspace} con inicio $0$, final $8$ y $300$ puntos.

    \item[(b)] Calcule los arreglos de posición $x(t)$ y $y(t)$ usando las ecuaciones cinemáticas entregadas.

    \item[(c)] Calcule los arreglos de velocidad $v_x(t)$ y $v_y(t)$. Recuerde que, aunque $v_x(t)$ sea constante, debe representarse como un arreglo con la misma dimensión que $t$ para poder graficarlo correctamente.

    \item[(d)] Encuentre el índice y el valor máximo de cada arreglo de velocidad, es decir, de $v_x(t)$ y $v_y(t)$. Además, calcule el promedio de ambos arreglos de velocidad.

    \item[(e)] Calcule los arreglos de aceleración $a_x(t)$ y $a_y(t)$. Para $a_x(t)$ puede usar un arreglo de ceros con la misma dimensión que el tiempo, por ejemplo mediante \texttt{np.zeros\_like(t)}. Para $a_y(t)$ debe construir un arreglo constante igual a $-g$.

    \item[(f)] Realice una figura con subgráficos usando \texttt{plt.subplots(3, 2)}. En la primera columna debe graficar las cantidades asociadas al movimiento vertical:
    \[
    y(t) \ \text{vs.} \ t,
    \qquad
    v_y(t) \ \text{vs.} \ t,
    \qquad
    a_y(t) \ \text{vs.} \ t.
    \]
    En la segunda columna debe graficar las cantidades asociadas al movimiento horizontal:
    \[
    x(t) \ \text{vs.} \ t,
    \qquad
    v_x(t) \ \text{vs.} \ t,
    \qquad
    a_x(t) \ \text{vs.} \ t.
    \]

    \item[(g)] Agregue títulos, nombres de ejes y una presentación limpia a cada gráfico, de modo que la figura permita comparar claramente la evolución horizontal y vertical del proyectil.
\end{itemize}
\end{tcolorbox}

\begin{solution}
Una posible solución consiste en construir primero el arreglo temporal, luego calcular las cantidades cinemáticas como arreglos de NumPy y finalmente graficarlas en una figura de $3 \times 2$ subgráficos.

\begin{pythonjupyter}
import numpy as np
import matplotlib.pyplot as plt

# Constantes físicas
y0 = 50          # altura inicial en metros
v0 = 35          # rapidez inicial en m/s
theta_deg = 55   # ángulo en grados
g = 9.8          # aceleración de gravedad en m/s^2

# Conversión del ángulo a radianes
theta = np.deg2rad(theta_deg)

# Arreglo temporal
t = np.linspace(0, 8, 300)

# Posiciones
x = v0 * np.cos(theta) * t
y = y0 + v0 * np.sin(theta) * t - 0.5 * g * t**2

# Velocidades
vx = v0 * np.cos(theta) * np.ones_like(t)
vy = v0 * np.sin(theta) - g * t

# Índices de velocidad máxima
idx_vx_max = np.argmax(vx)
idx_vy_max = np.argmax(vy)

# Valores máximos
vx_max = vx[idx_vx_max]
vy_max = vy[idx_vy_max]

# Promedios
vx_prom = np.mean(vx)
vy_prom = np.mean(vy)

print("Velocidad máxima en x:")
print("Índice:", idx_vx_max)
print("Valor:", vx_max)

print("\nVelocidad máxima en y:")
print("Índice:", idx_vy_max)
print("Valor:", vy_max)

print("\nPromedios:")
print("Promedio de vx:", vx_prom)
print("Promedio de vy:", vy_prom)

# Aceleraciones
ax = np.zeros_like(t)
ay = -g * np.ones_like(t)

# Figura con subgráficos
fig, axs = plt.subplots(3, 2, figsize=(10, 9))

# Columna izquierda: movimiento vertical
axs[0, 0].plot(t, y)
axs[0, 0].set_title("Posición vertical")
axs[0, 0].set_xlabel("t [s]")
axs[0, 0].set_ylabel("y(t) [m]")
axs[0, 0].grid(True)

axs[1, 0].plot(t, vy)
axs[1, 0].set_title("Velocidad vertical")
axs[1, 0].set_xlabel("t [s]")
axs[1, 0].set_ylabel("vy(t) [m/s]")
axs[1, 0].grid(True)

axs[2, 0].plot(t, ay)
axs[2, 0].set_title("Aceleración vertical")
axs[2, 0].set_xlabel("t [s]")
axs[2, 0].set_ylabel("ay(t) [m/s²]")
axs[2, 0].grid(True)

# Columna derecha: movimiento horizontal
axs[0, 1].plot(t, x)
axs[0, 1].set_title("Posición horizontal")
axs[0, 1].set_xlabel("t [s]")
axs[0, 1].set_ylabel("x(t) [m]")
axs[0, 1].grid(True)

axs[1, 1].plot(t, vx)
axs[1, 1].set_title("Velocidad horizontal")
axs[1, 1].set_xlabel("t [s]")
axs[1, 1].set_ylabel("vx(t) [m/s]")
axs[1, 1].grid(True)

axs[2, 1].plot(t, ax)
axs[2, 1].set_title("Aceleración horizontal")
axs[2, 1].set_xlabel("t [s]")
axs[2, 1].set_ylabel("ax(t) [m/s²]")
axs[2, 1].grid(True)

plt.tight_layout()
plt.show()
\end{pythonjupyter}
\end{solution}
